%% --------------------------------------------------------
\section{До визначення флуктуації числа частинок}\label{add:B}
%% --------------------------------------------------------

Розглянемо дві системи, що знаходяться у дифузному та тепловому контакті. При цьому одна із них, що будемо вважати термостатом, має об’єм $V_{\text{т}}$ значно більше ніж друга об’ємом $V_\text{т}$, яку будемо вважати системою (як на Рис.~\ref{pic:Pic14.6}). Отже $V / V_0 = p \ll 1$, де $V_0$ --- загальний об’єм системи і термостата.

Слідуючи тим же міркуванням, що і в \S~\ref{sec:MaxStatWeight} при оцінці флуктуації кількості частинок в двох рівних об’ємах, отримаємо, що імовірність спостерігати $N$ частинок в системі становить:

\begin{equation}
	P_N=\frac{N_0!}{N!\left(N_0-N\right)!}p^N\left(1-p\right)^{N_0-N}=C_{N_0}^Np^N\left(1-p\right)^{N_0-N}.
	\label{eq:AppB.1}
\end{equation}

Отриманий вираз є біномінальним розподілом. Використовуючи його можна знайти середнє і середній квадрат числа частинок $N$ в системі, проте очевидно, що середня кількість частинок, $\langle N \rangle  = N_0 p$, де $N_0$ --- повна кількість частинок в системі і термостаті. Для того, щоб знайти середній квадрат, скористаємся очевидним співідношенням $N^2 = N(N - 1) + N$. Тоді залишається знайти середнє значення $\langle N(N - 1) \rangle$:

\begin{multline}
\sum_{N=0}^{N_0} N(N-1)C_{N_0}^N p^N (1-p)^{N_0-N}
= \sum_{N=2}^{N_0} \frac{N_0!}{(N-2)! (N_0-N)!} p^N (1-p)^{N_0-N} \\
= N_0 (N_0 - 1)p^2 \sum_{N=2}^{N_0} \frac{(N_0 - 2)!}{(N-2)! (N_0-N)!} p^{N-2} (1-p)^{N_0-N} = \\
= N_0 (N_0 - 1)p^2 \sum_{M=0}^{N_0-2} \frac{(N_0 - 2)!}{(N-2)! (N_0 - 2-M)!} p^M (1-p)^{N_0-2-M} \\
= N_0 (N_0 - 1)p^2 \sum_{M=0}^{N_0-2} C_{N_0-2}^M p^M (1-p)^{N_0-2-M} = N_0 (N_0 - 1)p^2
	\label{eq:AppB.2}
\end{multline}

Звідси отримаємо, що середній квадрат числа частинок в системі становить:

\begin{equation}
	\langle N^2 \rangle = N_0 (N_0 - 1) p^2 + N_0 p = N_0^2 p^2 + N_0 p (1 - p) \cong N_0^2 p^2 + N_0 p = N_0^2 p^2 + \langle N \rangle
	\label{eq:AppB.3}
\end{equation}

де враховано, що $p \ll 1$. Остаточно, відносна флуктуація кількості частинок в системі становить:

\begin{equation}
	\frac{\langle \delta N \rangle}{\langle N \rangle} = \frac{\sqrt{\langle \Delta N^2 \rangle}}{\langle N \rangle} = \frac{\sqrt{(N_0^2 p^2 + \langle N \rangle) - N_0^2 p^2}}{\langle N \rangle} = \frac{1}{\sqrt{\langle N \rangle}}
	\label{eq:AppB.4}
\end{equation}